\subsection{Implementasi \textit{Client} Tenaga Kesehatan}

\textit{Client} tenaga kesehatan diimplementasikan menggunakan Tauri, SvelteKit, TypeScript, dan Rust. 
\textit{Frontend} SvelteKit digunakan untuk menampilkan halaman aplikasi, sedangkan \textit{backend} Rust pada Tauri digunakan untuk menyimpan kredensial lokal, melakukan operasi kriptografi, memanggil PRE, memanggil \textit{smart contract}, dan mendekripsi data yang diterima.

\textit{Client} ini digunakan oleh admin rumah sakit dan tenaga kesehatan. 
Admin rumah sakit dapat membuat \textit{activation key} untuk akun tenaga kesehatan. 
Tenaga kesehatan menggunakan \textit{activation key} tersebut untuk mengaktifkan akun, melengkapi profil, dan menghasilkan QR akses. 
QR akses berisi alamat IOTA tenaga kesehatan dan \textit{public key} PRE tenaga kesehatan. 
QR inilah yang kemudian dipindai oleh \textit{client} Android pasien ketika pasien memberikan akses.

\subsubsection{Halaman Aktivasi}

Halaman aktivasi digunakan untuk memasukkan \textit{activation key} dan data aktivasi akun rumah sakit atau tenaga kesehatan. 
Proses ini memastikan bahwa akun yang dibuat pada \textit{client} tenaga kesehatan berasal dari aktivasi yang sah.
Halaman aktivasi, halaman awal \textit{login}, dan halaman awal \textit{register} menggunakan tampilan yang sama, yaitu tampilan pada Gambar~\ref{halaman-login-hospital-pghd} pada Lampiran~\ref{app:tampilan-implementasi-client-hospital-pghd}.

\subsubsection{Halaman \textit{Register}}

Halaman \textit{register} digunakan untuk membuat akun lokal pada \textit{client} tenaga kesehatan setelah proses aktivasi berhasil. 
Pada halaman ini, \textit{client} menyiapkan kredensial dan kredensial yang diperlukan agar akun dapat digunakan pada aplikasi.
Halaman \textit{register} \textit{client} tenaga kesehatan dapat dilihat pada Gambar~\ref{halaman-register-hospital-pghd} pada Lampiran~\ref{app:tampilan-implementasi-client-hospital-pghd}.

\subsubsection{Halaman \textit{Login}}

Halaman \textit{login} digunakan untuk membuka akun lokal pada \textit{client} tenaga kesehatan. 
Setelah \textit{login} berhasil, \textit{client} dapat memuat profil, akses pasien, dan fitur lain sesuai peran pengguna.
Halaman \textit{login} \textit{client} tenaga kesehatan dapat dilihat pada Gambar~\ref{halaman-login-hospital-pghd} pada Lampiran~\ref{app:tampilan-implementasi-client-hospital-pghd}.

\subsubsection{Halaman Pelengkapan Profil}

Halaman pelengkapan profil digunakan untuk mengisi nama dan metadata profil tenaga kesehatan. 
Data profil ini digunakan ketika tenaga kesehatan menampilkan identitas dan QR akses kepada pasien.
Halaman melengkapi profil \textit{client} tenaga kesehatan dapat dilihat pada Gambar~\ref{halaman-complete-profile-hospital-pghd} pada Lampiran~\ref{app:tampilan-implementasi-client-hospital-pghd}.

\subsubsection{Halaman Profil}

Halaman profil digunakan untuk menampilkan data administratif tenaga kesehatan dan QR akses. 
QR tersebut memuat alamat IOTA tenaga kesehatan dan \textit{public key} PRE yang akan dipindai oleh pasien saat memberikan akses.
Halaman profil \textit{client} tenaga kesehatan dapat dilihat pada Gambar~\ref{halaman-profil-hospital-pghd} pada Lampiran~\ref{app:tampilan-implementasi-client-hospital-pghd}.

\subsubsection{Halaman \textit{Dashboard}}

Halaman \textit{dashboard} digunakan untuk melihat daftar akses yang diterima dari pasien. 
Dari halaman ini, tenaga kesehatan dapat memilih pasien dan jenis data yang akan dibuka sesuai \textit{capability} yang masih aktif.
Halaman \textit{dashboard} \textit{client} tenaga kesehatan dapat dilihat pada Gambar~\ref{halaman-dashboard-hospital-pghd} pada Lampiran~\ref{app:tampilan-implementasi-client-hospital-pghd}.

\subsubsection{Halaman PGHD}

Halaman PGHD digunakan untuk melihat daftar \textit{batch} PGHD pasien, membuka detail \textit{batch} PGHD, melihat \textit{provenance} data, dan menginvalidasi data PGHD apabila data terbukti tidak valid. 
Dekripsi dan verifikasi akhir tetap dilakukan pada \textit{client} tenaga kesehatan.
Halaman daftar PGHD \textit{client} tenaga kesehatan dapat dilihat pada Gambar~\ref{halaman-list-pghd-hospital} pada Lampiran~\ref{app:tampilan-implementasi-client-hospital-pghd}.

\subsubsection{Akses dan Dekripsi PGHD}

Setelah pasien memberikan akses PGHD, tenaga kesehatan memperoleh token akses dari PRE\@. 
Token tersebut digunakan untuk memanggil \textit{endpoint} daftar PGHD dan endpoint detail PGHD pada PRE, yang masing-masing dijelaskan pada Subbab~\ref{subsubsec:pre-endpoint-pghd-list} dan Subbab~\ref{subsubsec:pre-endpoint-pghd-detail}. 
\textit{Client} tenaga kesehatan tidak langsung memperoleh \textit{plaintext} PGHD dari PRE\@. 
PRE mengirimkan \textit{payload} terenkripsi, metadata \textit{on-chain}, \textit{cfrag}, \textit{capsule}, \textit{public key} PGHD pasien, dan material lain yang diperlukan untuk dekripsi. 
Dekripsi akhir tetap dilakukan pada \textit{client} tenaga kesehatan.
Halaman akses detail PGHD \textit{client} tenaga kesehatan dapat dilihat pada Gambar~\ref{halaman-detail-pghd-hospital} pada Lampiran~\ref{app:tampilan-implementasi-client-hospital-pghd}.

Pada saat membuka detail PGHD, \textit{client} melakukan beberapa tahap validasi. 
Pertama, \textit{client} membandingkan hash dari ciphertext yang diunduh dengan hash yang tercatat pada metadata. 
Kedua, \textit{client} memverifikasi \texttt{signature} atas \texttt{h\_cipher} menggunakan \textit{public key} PGHD pasien. 
Ketiga, \textit{client} melakukan proses \textit{re-encryption} menggunakan \textit{cfrag} dan \textit{private key} tenaga kesehatan. 
Keempat, \textit{client} melakukan dekripsi AES-GCM; tag autentikasi AES-GCM memastikan \textit{ciphertext}, kunci, \textit{nonce}, dan data tambahan yang dipakai pada enkripsi konsisten. 
Jika verifikasi hash atau \textit{signature}, maupun dekripsi AES-GCM, tidak valid, \textit{client} menampilkan peringatan kepada tenaga kesehatan dan \textit{memanggil endpoint} invalidasi PGHD pada PRE yang dijelaskan pada Subbab~\ref{subsubsec:pre-endpoint-pghd-invalidate}, dengan alasan invalid.

\subsubsection{Pemisahan PGHD dan \textit{Medical Record}}

PGHD dan \textit{medical record} ditampilkan sebagai dua jenis data yang berbeda. 
\textit{Medical record} tetap mengikuti mekanisme DecMed lama yang berorientasi pada data administratif dan klinis. 
PGHD menggunakan \textit{endpoint} dan token akses terpisah dengan tujuan akses \texttt{READ\_PGHD}. 
Token tersebut dibuat melalui \textit{endpoint key} PRE pada Subbab~\ref{subsubsec:pre-endpoint-keys}. 
Pemisahan ini membuat tenaga kesehatan dapat memiliki akses ke PGHD tanpa harus memiliki akses update \textit{medical record}, atau sebaliknya sesuai akses yang diberikan pasien.

\subsubsection{\textit{Command} Tauri yang Digunakan}

\textit{Backend} Tauri menyediakan command yang dipanggil oleh \textit{frontend} dan dapat dilihat pada Tabel~\ref{tab:implementasi-command-hospital-client}.

\begin{styledxltabular}{|>{\raggedright\arraybackslash}p{0.32\textwidth}|>{\raggedright\arraybackslash}X|}
    \hiderowcolors
    \caption{\textit{Command} pada \textit{Client} Tenaga Kesehatan}\label{tab:implementasi-command-hospital-client} \\
    \hline
    \rowcolor{headerBg}
    \textbf{Command} & \textbf{Fungsi} \\
    \hline
    \showrowcolors
    \endfirsthead
    \hline
    \rowcolor{headerBg}
    \textbf{Command} & \textbf{Fungsi} \\
    \hline
    \endhead
    \hline
    \endfoot
    \hline
    \endlastfoot
    \texttt{get\_profile} & Mengambil profil tenaga kesehatan dan metadata yang diperlukan untuk menampilkan QR code. \\
    \hline
    \texttt{get\_medical\_record} & Membaca \textit{medical record} pasien berdasarkan token akses baca \textit{medical record}. \\
    \hline
    \texttt{new\_medical\_record} & Membuat \textit{medical record} baru apabila tenaga kesehatan memiliki akses update. \\
    \hline
    \texttt{update\_medical\_record} & Memperbarui \textit{medical record} pasien berdasarkan akses update yang diberikan pasien. \\
    \hline
    \texttt{get\_pghd\_list} & Mengambil daftar metadata \textit{batch} PGHD pasien melalui PRE\@. \\
    \hline
    \texttt{get\_pghd} & Mengambil, melakukan \textit{re-encryption}, mendekripsi, dan memverifikasi detail \textit{batch} PGHD\@. \\
    \hline
    \texttt{invalidate\_pghd} & Mengirim permintaan invalidasi PGHD ketika data terbukti rusak atau tidak sah. \\
    \hline
\end{styledxltabular}
